\documentclass[12pt,a4paper,landscape]{article}
\usepackage[T1]{fontenc}
\usepackage[utf8]{inputenc}
\usepackage[ngerman]{babel}
\usepackage[margin=14mm]{geometry}
\usepackage{amsmath,amssymb}
\usepackage{tikz}
\usetikzlibrary{arrows.meta,calc,patterns}
\pagestyle{empty}
\setlength{\parindent}{0pt}
\renewcommand{\normalsize}{\fontsize{16}{20}\selectfont}
\renewcommand{\large}{\fontsize{18}{22}\selectfont}
\renewcommand{\Large}{\fontsize{20}{25}\selectfont}
\renewcommand{\LARGE}{\fontsize{24}{29}\selectfont}
\renewcommand{\small}{\fontsize{13}{16}\selectfont}
\definecolor{petrol}{RGB}{0,90,91}
\definecolor{orange}{RGB}{166,63,30}
\definecolor{blue}{RGB}{33,70,130}
\tikzset{every picture/.style={line width=0.9pt}, >={Stealth[length=2.7mm]},
  lab/.style={font=\large}, smalllab/.style={font=\normalsize},
  axis/.style={->,line width=0.9pt}, helper/.style={dashed,gray}}
\newcommand{\kopf}[2]{\noindent{\LARGE\bfseries #1}\hfill{\large #2}\par\vspace{2mm}\hrule\vspace{4mm}}
\newcommand{\folge}[1]{\vfill\noindent\rule{\textwidth}{0.4pt}\par\vspace{2mm}{\small\textbf{Live aufbauen:} #1}}
\begin{document}
\normalsize

\kopf{1 · Ein Bild aus Pixeln}{Einstieg}
\begin{minipage}[t]{.46\textwidth}
\vspace{0pt}\begin{tikzpicture}[scale=.72]
  \fill[petrol!20] (3,2) rectangle (4,3);
  \draw[step=1,gray] (0,0) grid (7,5);
  \draw[very thick] (0,0) rectangle (7,5);
  \node[lab] at (3.5,2.5) {$\bullet$};
  \draw[->,orange,thick] (5,4.3) -- (3.8,2.8);
  \node[lab,anchor=west] at (5,4.4) {aktuelles Pixel};
  \draw[axis] (0,-.5) -- (7.4,-.5) node[right] {$x$};
  \draw[axis] (-.5,0) -- (-.5,5.4) node[above] {$y$};
  \node[smalllab,anchor=north east] at (0,0) {$(0,0)$};
\end{tikzpicture}
\end{minipage}\hfill
\begin{minipage}[t]{.50\textwidth}\vspace{0pt}
\Large
Für jedes Pixel wird die Farbe berechnet:
\[\boxed{\text{Farbe}=f(x,y,t)}\]
\normalsize
\begin{tabular}{ll}
\texttt{fragCoord} & Position des aktuellen Pixels \\
\texttt{iResolution.xy} & Bildbreite und Bildhöhe \\
\texttt{iTime} & Zeit in Sekunden
\end{tabular}
\par\medskip
\texttt{mainImage} berechnet eine Farbe pro Pixel.
\end{minipage}
\folge{Raster zeichnen → ein Pixel markieren → Position, Bildgröße und Zeit benennen → Farbformel.}
\newpage

\kopf{2 · Farbe und Zeit}{Demo 1}
\begin{minipage}[t]{.47\textwidth}\vspace{0pt}
\begin{tikzpicture}[x=1cm,y=1cm]
 \foreach \y/\name/\shade in {3/Rot/orange,2/Grün/petrol,1/Blau/blue}{
   \node[lab,anchor=east] at (0,\y) {\name};
   \draw (0.3,\y-.22) rectangle (5.5,\y+.22);
   \fill[\shade!70] (0.3,\y-.22) rectangle (3.6,\y+.22);
   \node[smalllab] at (0.3,\y-.6) {0}; \node[smalllab] at (5.5,\y-.6) {1};
 }
\end{tikzpicture}
\end{minipage}\hfill
\begin{minipage}[t]{.50\textwidth}\vspace{0pt}
\Large\[\text{Farbe}=(R,G,B)\]
\normalsize Jeder Wert liegt zwischen 0 und 1.\\
$(0,0,0)$ ist Schwarz, $(1,1,1)$ ist Weiß.
\par\medskip
Mit Zeit ändert sich die Farbe überall zugleich:
\[R(t)=0{,}5+0{,}5\sin(t)\]
Grün und Blau schwingen zeitversetzt.
\end{minipage}
\folge{RGB-Anteile eintragen → zwei Farben vergleichen → Zeit als Eingabe ergänzen → Wertebereich des Sinus erklären.}
\newpage

\kopf{3 · Position wird Farbe}{Demo 2}
\begin{minipage}[t]{.49\textwidth}\vspace{0pt}
\begin{tikzpicture}[x=1cm,y=1cm]
 \draw[thick] (0,0) rectangle (7,4);
 \draw[axis] (0,0) -- (7.5,0) node[right] {$x$};
 \draw[axis] (0,0) -- (0,4.4) node[above] {$y$};
 \fill[orange] (3.5,2.5) circle (2.5pt);
 \draw[helper] (3.5,0) -- (3.5,2.5) -- (0,2.5);
 \node[lab,above right] at (3.5,2.5) {$(x,y)$};
 \node[smalllab,below] at (3.5,0) {$x$};
 \node[smalllab,left] at (0,2.5) {$y$};
 \node[smalllab,below] at (7,0) {Breite $W$};
 \node[smalllab,left] at (0,4) {Höhe $H$};
\end{tikzpicture}
\end{minipage}\hfill
\begin{minipage}[t]{.48\textwidth}\vspace{0pt}
\Large\[u=\frac{x}{W},\qquad v=\frac{y}{H}\]
\normalsize Positionsfarbe: $(R,G,B)=(u,v,0)$.\par\medskip
Zwei Farben entlang der x-Achse mischen:
\[\text{Farbe}=(1-u)\,\text{links}+u\,\text{rechts}\]
$u=0$: links; \quad $u=1$: rechts.
\par\medskip
Getrennt normierte Achsen sind für Abstände noch nicht gleich skaliert.
\end{minipage}
\folge{Pixel und Auflösung markieren → normieren → Koordinaten als Farben zeigen → zwei Farben mischen.}
\newpage

\kopf{4 · Vom Abstand zum Kreis}{Demo 3}
\begin{minipage}[t]{.48\textwidth}\vspace{0pt}
\begin{tikzpicture}[x=1cm,y=1cm]
 \draw[axis] (-3.4,0) -- (3.5,0) node[right] {$p_x$};
 \draw[axis] (0,-2.4) -- (0,2.6) node[above] {$p_y$};
 \draw[petrol,very thick] (0,0) circle (2);
 \fill (0,0) circle (2pt) node[below left] {$M$};
 \fill[orange] (1.45,1.05) circle (2.5pt) node[above right] {$P$};
 \draw[orange,thick,->] (0,0) -- (1.45,1.05) node[midway,above left] {$\vec p$};
 \draw[helper] (1.45,0) -- (1.45,1.05) -- (0,1.05);
 \draw[petrol,->] (0,0) -- (-1.4,1.43) node[midway,above left] {$r$};
\end{tikzpicture}
\end{minipage}\hfill
\begin{minipage}[t]{.49\textwidth}\vspace{0pt}
Bildmitte als Ursprung, durch Bildhöhe skalieren:
\[\vec p=\frac{\texttt{fragCoord}-\tfrac12\,\texttt{iResolution.xy}}{\texttt{iResolution.y}}\]
Gleiche Einheiten auf beiden Achsen halten den Kreis rund.
\[d^2=p_x^2+p_y^2=\operatorname{dot}(\vec p,\vec p)\]
\[\text{innen: }d^2<r^2\qquad\text{außen: }d^2\ge r^2\]
\end{minipage}
\folge{Mittelpunkt und Pixel einzeichnen → Koordinaten zentrieren → Abstand → Radiusvergleich und zwei Farben.}
\newpage

\kopf{5 · Von Kreisen zu Metaballs}{Demo 3}
\begin{minipage}[t]{.48\textwidth}\vspace{0pt}
\begin{tikzpicture}[x=1cm,y=1cm]
 \draw[petrol,very thick] (-1.45,1.7) circle (1.25);
 \draw[blue,very thick] (1.45,1.7) circle (1.25);
 \fill[petrol] (-1.45,1.7) circle (3pt) node[below] {$M_1$};
 \fill[blue] (1.45,1.7) circle (3pt) node[below] {$M_2$};
 \node[lab] at (0,-.15) {Zwei Felder beeinflussen jeden Punkt.};
 \draw[->,thick] (-2,-.6) -- (2,-.6);
 \node[smalllab] at (0,-1.05) {Bewegung der Mittelpunkte mit $t$};
\end{tikzpicture}
\end{minipage}\hfill
\begin{minipage}[t]{.49\textwidth}\vspace{0pt}
Ein Feld wird mit wachsendem Abstand kleiner:
\[E_i(P)=\frac{a}{\lVert P-M_i\rVert^2+\varepsilon}\]
Die Felder addieren sich:
\[E(P)=E_1(P)+E_2(P)\]
\[E(P)>\text{Schwelle}\quad\Rightarrow\quad\text{Vordergrund}\]
Zwischen den Zentren kann die Summe die Schwelle überschreiten: Die Formen verbinden sich. $\varepsilon$ verhindert Division durch null.
\end{minipage}
\folge{Zwei Kreise → abfallende Felder → Werte an einem Punkt addieren → Schwelle → Bewegung.}
\newpage

\kopf{6 · Weiche Übergänge}{Polynom und Kante}
\begin{minipage}[t]{.48\textwidth}\vspace{0pt}
\begin{tikzpicture}[x=5cm,y=3.5cm]
 \draw[axis] (-.08,0) -- (1.18,0) node[right] {$t$};
 \draw[axis] (0,-.08) -- (0,1.12) node[above] {$s$};
 \draw[gray,dashed] (0,0) -- (1,1) node[above right,font=\small] {linear};
 \draw[petrol,very thick,domain=0:1,samples=60] plot (\x,{3*\x*\x-2*\x*\x*\x});
 \fill (0,0) circle (.55pt) (1,1) circle (.55pt);
 \node[below] at (1,0) {1}; \node[left] at (0,1) {1};
 \node[petrol,anchor=west] at (.43,.23) {$3t^2-2t^3$};
\end{tikzpicture}
\end{minipage}\hfill
\begin{minipage}[t]{.49\textwidth}\vspace{0pt}
Im Übergang läuft der Mischparameter $t$ von 0 bis 1.
\[s(t)=at^3+bt^2+ct+d\]
Die Bedingungen $s(0)=0$, $s(1)=1$ und $s'(0)=s'(1)=0$ ergeben $a=-2$, $b=3$, $c=d=0$:
\[s(t)=3t^2-2t^3\]
Die Kurve startet und endet waagerecht.
\[\text{Farbe}=(1-s)\,\text{innen}+s\,\text{außen}\]
Anwendung: Kante eines Kreises oder einer Feldschwelle.
\end{minipage}
\folge{Harte Kante → Übergangsintervall → lineare Mischung → Randbedingungen → kubisches Polynom.}
\newpage

\kopf{7 · Sinus: Form und Bewegung}{Wellen}
\begin{minipage}[t]{.49\textwidth}\vspace{0pt}
\begin{tikzpicture}[x=.9cm,y=1cm]
 \draw[axis] (-.2,0) -- (7.2,0) node[right] {$x$};
 \draw[axis] (0,-1.7) -- (0,1.8) node[above] {Höhe};
 \draw[petrol,very thick,domain=0:6.8,samples=100] plot (\x,{1.1*sin(\x*90)});
 \draw[blue,dashed,thick,domain=0:6.8,samples=100] plot (\x,{1.1*sin(\x*90-60)});
 \draw[<->,orange,thick] (.7,0) -- (.7,1.1) node[midway,right] {$A$};
 \node[smalllab] at (3.4,-2.15) {durchgezogen: früher; gestrichelt: später};
\end{tikzpicture}
\end{minipage}\hfill
\begin{minipage}[t]{.49\textwidth}\vspace{0pt}
\[h(x,t)=A\sin(kx-\omega t)\]
\begin{tabular}{ll}
$A$ & Amplitude: wie hoch? \\
$k$ & räumliche Frequenz: wie eng? \\
$\omega t$ & Phase: wie verschoben?
\end{tabular}
\par\medskip
Ein mögliches Motiv: Wasser unter der Linie, Himmel darüber. Die Kurve dient hier zur Erklärung.
\end{minipage}
\folge{Wiederholung der Farbschwingung → Höhe über x → Amplitude und Frequenz → Zeit verschiebt die Phase.}
\newpage

\kopf{8 · Sinus durch ein Polynom annähern}{Taylor}
\begin{minipage}[t]{.49\textwidth}\vspace{0pt}
\begin{tikzpicture}[x=.72cm,y=1.45cm]
 \draw[axis] (-4.1,0) -- (4.2,0) node[right] {$x$};
 \draw[axis] (0,-1.35) -- (0,1.45);
 \draw[petrol,very thick,domain=-3.7:3.7,samples=120] plot (\x,{sin(\x r)});
 \draw[orange,dashed,thick,domain=-3:3,samples=120] plot (\x,{\x-\x*\x*\x/6+\x*\x*\x*\x*\x/120});
 \node[petrol,anchor=west] at (2.1,1.15) {$\sin(x)$};
 \node[orange,anchor=west] at (2.1,.85) {Polynom};
\end{tikzpicture}
\end{minipage}\hfill
\begin{minipage}[t]{.49\textwidth}\vspace{0pt}
\[\sin(x)\approx x-\frac{x^3}{6}+\frac{x^5}{120}\]
\textbf{Entwicklungspunkt:} $x=0$.
\par\medskip
In der Nähe von 0 stimmen die Kurven gut überein. Weiter entfernt wächst die Abweichung.
\par\medskip
Die Approximation ist ein Polynom, der Sinus eine periodische Funktion.
\end{minipage}
\folge{Sinus als Referenz zeichnen → Polynomnäherung ergänzen → Nähe bei 0 und Abweichung besprechen.}
\newpage

\kopf{9 · Value Noise: Zufall wird weich}{Noise}
\begin{minipage}[t]{.49\textwidth}\vspace{0pt}
\begin{tikzpicture}[scale=1.25]
 \draw[gray] (0,0) grid (3,3);
 \fill[petrol!15] (1,1) rectangle (2,2);
 \foreach \x/\y/\s in {1/1/a,2/1/b,1/2/c,2/2/d}{\fill[petrol] (\x,\y) circle (2pt); \node[above right,font=\small] at (\x,\y) {$\s$};}
 \fill[orange] (1.55,1.35) circle (2.2pt) node[below right] {$P$};
 \draw[orange,->] (1.55,1.35) -- (2.5,.45);
 \node[smalllab,anchor=west] at (2.5,.4) {Position in der Zelle};
\end{tikzpicture}
\end{minipage}\hfill
\begin{minipage}[t]{.49\textwidth}\vspace{0pt}
\begin{enumerate}\setlength{\itemsep}{5pt}
 \item Jedem Gitterpunkt einen reproduzierbaren Pseudozufallswert zuordnen.
 \item Innerhalb der Zelle zwischen den vier Eckwerten mischen.
 \item Den Mischfaktor mit $s(t)=3t^2-2t^3$ glätten.
\end{enumerate}
Ergebnis: ein zusammenhängendes, weich veränderliches Muster.
\end{minipage}
\folge{Gitter und Eckwerte → Punkt in einer Zelle → erst horizontal, dann vertikal mischen → glätten.}
\newpage

\kopf{10 · fBm: Grob und fein überlagern}{fBm}
\begin{center}
\begin{tikzpicture}[x=1cm,y=1cm]
 \foreach \x/\step/\label in {0/1/Grob,5/.5/Mittel,10/.25/Fein}{
   \draw[step=\step,gray] (\x,0) grid +(3,3);
   \draw[very thick] (\x,0) rectangle +(3,3);
   \node[lab] at (\x+1.5,-.45) {\label};
 }
 \node[font=\Huge] at (4,1.5) {$+$};
 \node[font=\Huge] at (9,1.5) {$+$};
 \draw[->,thick] (13.4,1.5) -- (14.4,1.5);
 \node[lab,align=left,anchor=west] at (14.5,1.5) {Wolken-\\struktur};
\end{tikzpicture}
\end{center}
\[F(p)=\frac{0{,}5\,N(p)+0{,}25\,N(2p)+0{,}125\,N(4p)+0{,}0625\,N(8p)}{0{,}9375}\]
\begin{center}$N$: Value Noise. Das Bild zeigt drei Maßstäbe; die Formel verwendet vier Oktaven. Jede verdoppelt die Frequenz und halbiert das Gewicht.\end{center}
\folge{Ein Noise-Muster → kleinere Strukturen → schwächer gewichtet addieren → Gewichtssumme erklären.}
\end{document}
